\subsection{Seventh letter}

\begin{quotationx}
A few quick points about this letter. Masson-Oursel was mentioned in Letter III\footnote{\url{https://gornahoor.net/?p=4739}}. Apparently he did an about-face and severely criticizes Guenon. Guenon's response is instructive; it shows how to deal with criticism.

Guenon and de Giorgio clearly have a wide range of interests. The topic of Virgil is interesting, since Guenon indicates there is deeper meaning to him. Although Guenon never did write about Virgil, De Giorgio did through his ideas of the two Romes and the influence of Dante. I don't know which book by Wallace he is referring to.

At this early date, we can see that De Giorgio is pushing Guenon into the topic of Buddhism. Regarding the Sufi group in question, please see the comment by Saladin below.

It is curious that Guenon would not attend a tariqa that accommodated Europeans. So much for tying him to European racism or white nationalism. He claims their presence would distort the teachings. We see something similar today with all the Tibetan Buddhist centers in the West, in which the tenets of Buddhism are distorted if not ignored altogether.
\end{quotationx}


Paris, 31 December 1927

I received your two letters and I had wanted to respond to them sooner; but I have less and less free time, so that I've had to postpone until now, all the more reason that I intend to write you at length. I believe that, even today, I am going to be obliged to respond only to part of what you tell me and to put off the rest until another time.

First of all, Evola wrote to me twice recently and, according to what you tell me, I see what he tells me is pretty much the same that he wrote to you. I ask myself what can be that ``action'' that he counts on to lead with his journal and book, not only in Italy but also in neighboring countries; we will certainly see in the end.

He also sent me an article titled ``Anti-philosophical Fascism and Mediterranean Tradition'' that he published in the ``Critica Fascista''; perhaps you have seen it. There are some correct things, especially in the first part, but then his ``anti-Christianity'' reappears so that the editors of the journal had to add a note making some reservations on this point. On the other hand, I do not see very clearly what he means by the ``Mediterranean Tradition'', and I fear that, in the idea that he makes of it, there was a certain element of fantasy. I haven't had the time to respond to him. It is necessary all the same that I do so one of these days. As to the question of the collaboration with ``Ur'', I think just like you, so I am no more inclined than you to accept the limits that he wants to impose on us. Moreover, as he sees ``polemics'' in those things where my intention was quite otherwise, I don't know what I could do in order to remove those limits, which correspond only to a more questionable evaluation.

In the last number of \emph{Voile d'Isis}, there is a short note on ``Ur'', in which it is said that that journal is dedicated to the study of psychology! I don't know who wrote that, but it is certainly someone who has not read a single one of its articles. It is likely that the expression ``science of the I'' led to that misunderstanding. If Evola saw that, he must not have been very satisfied.

Two or three days ago, I received a card from Reghini who had not given me any sign of life for close to a year. He told me that he is going to write me before long and that he will explain the reasons for his silence.

My new book arrived two weeks ago, although it was due to appear in November. It would have been quite astonishing had the printer not been so late, because that is their habit. Also it was decided to publish the volume in January after the holidays because this period of gift giving is totally unfavorable.

We saw Hacking recently and we told him that you were in Blois for vacation. He regretted that you could not stop by Paris because he would have been very happy to see you again.

Masson-Oursel gave in the ``Revue Critique'' (I do not know too much about that publication, having received only a clipping) a short review of the \emph{King of the World} that I will have to quote it literally for you since it is worth the trouble:

\begin{quotex}
Leibniz loved to say that there is gold in the manure of scholasticism; he probably finds in it --- even among the alchemists --- that universal symbolic that the Gnostics, Hindus, Chinese, Kabbalists had advocated. Unfortunately, R. Guenon does not try to extract that gold; only criticism could claim that. He takes everything in good faith, provided that it is a traditional fact and he does not doubt that everything corresponds to everything; by that he proves himself to be in the line of the symbolists. He possesses the information but he accepts it does not matter what. Criticism would be to his eyes a poor enterprise which would discredit the seeker and be superfluous for an author who believes he has the metaphysical truth! We know him enough by his other works. What he takes for strength and lucidity, to our eyes impairs the value of too extensive knowledge. We are frightened that in a hundred small tidy pages, he claims to reveal the ultimate knowledge on the swastika, Aum and Manu, the luz and the Shekinah, the Graal, the Mages, and the Old Man of the Mountain, and enigmas without number. Even if he divines something correct here or there, what does the result matter without demonstration? And the proof that there is only one symbolic among the diversity of religions and philosophies?
\end{quotex}


What do you say about that? What is most stupefying is that he assured me a few months ago that he was coming around more and more to my point of view! It seems that he is on the contrary further from it than ever; he is obviously incapable of understanding. Nothing counts for him outside of ``criticism'' and analysis, that is, what I consider precisely as non-existent. He imagines, like all his colleagues, that volumes of heavy erudition are necessary to treat the least point of detail. If I had believed it useful to answer him, I would have been able to draw his attention that I am quite far from accepting no matter what information as he claims, since I refuse to hold the least account of that which comes from official orientalists. But that would have been perfectly useless. One cannot change the mentality of those people, and the best thing is to continue on one's way without worrying about what they say or think.

I haven't had the chance to take up Virgil for quite a while; it is certain that things are found in him that it is rather difficult to understand now, but I am convinced that, if one had the time, it would be useful to do some research into the point of view that interests us. In any case, the few interpretations that you indicate to me seem totally justified.

I was not able to find in the passages that you point out to me in P. Wallace's book, because I loaned it to someone who hasn't returned it to me yet. Thus, I will have to look again a little later and then I will speak to you about it again. For the reflection that I made to you another time on this subject and what you tell me, what you formulate corresponds actually very well to what I had meant.

I spoke to Bacot about the German translation of the \emph{Songs of Milarepa}. He knows it, but, not knowing German very well, he was unable to examine it too closely to see what it precisely meant. Moreover, it appears to be very incomplete. The passages that you cite for me do not seem very clear and it was lacking any notes to explain them; but is the translator capable of doing it? You are certainly correct re the central artery, which can only be \emph{sushumna} and it is more than probable that the ``wind'' is \emph{prana}. I ask myself even if, where it is a question of the ``five winds'', there would have been a connection to make with the five \emph{vayus} which are the modalities of \emph{prana}. The meanings of mudra are certainly what you indicate. As to Dharmadhatu, literally ``seed of the Law'', I know it is a designation of Buddha, but that does not clarify at all the meaning of the phrase where this word is found. You know that Vajra means at once ``lightning'' and ``diamond''; there would be a lot of research to do regarding that.

I haven't yet had the time to look for the text in the \emph{Bhagavad Gita}. I take note of the passage which you spoke to me about and will try to verify it one of these days.

Thanks in advance for the Vedantic poem \emph{Ashtavakra Gita} that you promised to send me; I think that you will not be too much in a hurry for me to return it to you, because at this moment I can't even read what I have. Thanks also for the information you give me on the topic of the texts in its translation into prose; but it will be difficult for me to find that because I never go into any bookstores.

The meaning of \emph{sakshi} is ``witness'' or ``observer''; I think that \emph{sakshipurusha} must be taken as an equivalent of the ``two birds in the tree'', one who ``looks without eating'', and who is actually the personality, while the other, who eats (that is, who is engaged in actions and its consequences), is jivatma or the individuality. The translation of the shloka couplet in the poem in question seems to be exact at first glance. When I have a moment, I will look again more closely.

The letter Kha symbolizes the atmosphere and even the sky, rather than air as an element. The designation of ``triform'' Agni can have several meanings; it is possible that, in the case here, it is related to the celestial fire, to the terrestrial fire and the vital fire. The three colours that you mention are more customarily related to the \emph{gunas} than to the elements: \emph{sattwa}, white: \emph{rajas}, red (it is even the meaning proper to the root rauf); \emph{tamas}, black (obscurity). For the 0 between 3 and 2 (in place of unity), what you say appears to me totally correct, or in any case very plausible as the interpretation. As to the passages of the commentary indicating the definition of the mukti following different schools, it is still a thing that I have to put off to the next time, not wanting to give you a translation that would risk being inexact.

I am happy that Taillard responded to you and that he was able to send you a photograph of Kheireddine. If you happen to make a copy as you hope, I would be very happy to have it. A peculiar thing, the same day that I received your letter, I got one from Jossot who also spoke to me about Kheireddine. He says that he is dead even though he arranged to instruct him; but I don't really know if Jossot would be able to profit much from his teachings. You will be very kind to tell me if you get a new response from Tallard. The Sheik of which he spoke to you is certainly Sheik Mostafa ben Alioua, the head of the Alawis. I have a summary of the teachings from that school, written by the Sheik Si Mohammed El Aïd's secretary, with a French translation made by Taillard, and a copy of which I could send you if that interests you. I believe you said that the Alawis now have a center in Paris, intended exclusively for Arabs and Kabyles. Although they have tried to put me in touch with than, I have not had the time for it, though it is really close to here. I will certainly have to go see it, because it would be more interesting than the branch which is being formed and into which Europeans will be admitted, as it seems to me you also said. The introduction of Western elements is too easily a cause for deviation.

Excuse my incomplete response; I will try to not wait so long to write to you.

I think that you are not in the mountains at the time of cold weather; haven't you suffered too much from it?

P.S. I almost forgot to thank you for your gracious gift of stamps.

\begin{quotex}
Update 24 September 2022.

H/T Saladin on 2012-09-25 at 21:04 said:

Dear Cologero,

The Alawi order mentioned by Guenon in the letter is not the same as the Alwai sect to which Assad family of Syria belongs. The former is a Sufi order whereas the latter is sect (one which many muslims believe is heretical) mainly predominant in western Asia (Syria, Kurdistan, etc.) Here is a link for the sufi order: \url{https://en.wikipedia.org/wiki/Ahmad\_al-Alawi}

This is a link for the sect: \url{https://en.wikipedia.org/wiki/Alawi}
\end{quotex}

\flrightit{Posted on 2012-09-24 by Cologero}

\begin{center}* * *\end{center}

\begin{footnotesize}
\begin{sffamily}
\texttt{Avery on 2012-09-25 at 01:27 said:}

``Nothing counts for him outside of `criticism' and analysis, that is, what I consider precisely as non-existent; he imagines, like all his colleagues, that volumes of heavy erudition are necessary to treat the least point of detail. ''

A damning indictment of the cult-crit impulse then and now, which one has to keep in mind whenever your reading passes from the metaphysical world into the modern academic one.

\hfill

\texttt{Saladin on 2012-09-25 at 21:04 said:}

Dear Cologero,

The Alawi order mentioned by Guenon in the letter is not the same as the Alwai sect to which Assad family of Syria belongs. The former is a Sufi order whereas the latter is sect (one which many muslims believe is heretical) mainly predominant in western Asia (Syria, Kurdistan, etc.) Here is a link for the sufi order: \url{http://en.wikipedia.org/wiki/Ahmad\_al-Alawi}

This is a link for the sect:
\url{https://en.wikipedia.org/wiki/Alawi}


\hfill

\texttt{Cologero on 2012-09-25 at 22:03 said:}

Thanks, Saladin, for pointing that out. I'll make the correction.

\hfill

\texttt{Cologero on 2012-09-26 at 00:17 said:}

Someone recommended to me the following authors as sources for an understanding of Virgil. Since I don't know that he wants his name to appear here, I shall leave it as anonymous.

A short addition on the commentators of Virgil which merit being read for highlighting the intentions of the Latin poet:

It is worth the effort to study the alchemical exegesis proposed by Petrus Bonus (XV century), Bracesco (XVI), d'Espagnet (XVII), Thomas Vaughan (XVII), Maïer (XVII), dei Conti (XVII), Pernety (XVIII), and more recently, by d'Hooghvorst (XX).


\hfill

\end{sffamily}
\end{footnotesize}

